\begin{lemma}
Assume the eventual parameter comparisons, \(n>n_0\), the rounded identities
\[
m=\noderef{TwoBiteNaturalReducedVertexCount}(n),\qquad
p=\frac12\sqrt{\frac{\log n}{n}},
\]
the fixed-set scale identity
\[
|I|=\noderef{TwoBiteNaturalIndependenceScale}(1+\varepsilon,n),
\]
and the fiber-degree event.  If a red staged-open coordinate
\(\operatorname{red}(q)\) has an opposite-color witness \(b\), meaning
\[
q\in\noderef{TwoBiteBasePairs}
\bigl(\pi_R^\omega(\noderef{TwoBiteX}(\omega,I,\operatorname{blue}(b)))\bigr),
\]
then it has such a witness \(b\) which is large, medium, or small.  Symmetrically,
if a blue staged-open coordinate \(\operatorname{blue}(q)\) has an
opposite-color witness \(r\), meaning
\[
q\in\noderef{TwoBiteBasePairs}
\bigl(\pi_B^\omega(\noderef{TwoBiteX}(\omega,I,\operatorname{red}(r)))\bigr),
\]
then it has such a witness \(r\) which is large, medium, or small.
\end{lemma}
\begin{proof}
We prove the red statement; the blue statement is identical with the colors
interchanged.  Let \(\operatorname{red}(q)\) be staged-open and suppose that
\(b\) is an opposite-color witness.  By
\(\noderef{TwoBiteClassSmallMediumLargeOrHuge}\), the blue base vertex \(b\)
is small, medium, large, or huge.  In the small, medium, and large cases, the
same witness \(b\) gives the desired large/medium/small conclusion.

It remains to exclude the huge case.  If \(b\) is huge, then the red half of
\(\noderef{FixedSetHugeOppositeWitnessTouched}\), applied with the present
eventual-comparisons, rounded-parameter, scale, and fiber-degree hypotheses,
says that \(\operatorname{red}(q)\) is pre-terminal recorded and in particular
is not in \(\noderef{TwoBiteStagedOpenPairs}(\omega,\varepsilon,I)\).  This
contradicts the staged-open assumption.  Therefore the huge alternative is
impossible, and the witness may be chosen large, medium, or small.  The blue
case uses the blue half of
\(\noderef{FixedSetHugeOppositeWitnessTouched}\) in the same way.
\end{proof}
